\chapter{The Same Blocks in Verilog}
\label{ch:bridge}

In Part~I you built circuits from blocks: you drew gates in the Designer,
assembled an adder, remembered an alarm in a flip-flop, and carried along the running
``Security Panel'' project. Now the same things are said in words, in the Verilog language.
The question of this chapter is: how is each familiar block written in code, and how can you
make sure that you got exactly the circuit you assembled by hand?

We will look for the answer in one and the same way. We take a block about which we already know
everything from Part~I, describe it in a few lines, click
``Synthesize'' and compare the RTL tab with our own circuit. The waveforms
confirm the truth table or the behavior over time. At the end of the chapter all the
blocks come together in the ``Security Panel'' in Verilog.

\begin{table}[H]
\centering
\caption{A block from Part~I and its Verilog form}
\small
\begin{tabularx}{\textwidth}{l X X}
\toprule
\textbf{Block} & \textbf{Circuit of blocks} & \textbf{Verilog line} \\
\midrule
Logic & AND, OR, NOT in the Designer & \code{assign alarm = door \& (night | armed);} \\
Adder & a chain of full adders & \code{assign \{co, s\} = a + b + ci;} \\
Comparison & the ``A greater than B'' comparator & \code{assign loud = noise > limit;} \\
Decoder & four AND gates & \code{assign y = en ? (4'd1 << a) : 4'd0;} \\
Flip-flop & a D flip-flop with a clock & \code{always @(posedge clk) q <= d;} \\
Counter & a register and a ``+1'' adder & \code{q <= q + 1;} \\
State machine & a state register and transition logic & \code{case (state) ... endcase} \\
\bottomrule
\end{tabularx}
\end{table}

All the source files are in the folder \code{code/00\_theory/}. Open any file in
Logra (the ``From scratch'' project, the ``Code'' tab, the ``Create'' group), click
``Synthesize (Ctrl+Enter)'' and follow along with the book. Look at the results in
the tabs of the ``Check'' group: RTL, Waveforms and Lab.

\section{The Adder: Formulas Turn into Gates}

The full adder from the chapter ``Arithmetic Blocks'' is described by two formulas.
In Verilog each formula is one \code{assign} line.

\vfile{\codepath/00_theory/full_adder.v}{code/00\_theory/full\_adder.v}

After synthesis the RTL tab shows exactly what we drew in Part~I:
two XOR elements, two ANDs and an OR (Fig.~\ref{fig:b-adder-rtl}). Each
gate is labeled with a line number in the source text, and the wires carry the names of
the signals.

\shot{br-adder-rtl}{The RTL diagram of a full adder: the gates from the textbook, with labels taken from the code}{fig:b-adder-rtl}

Let us check the truth table. The ``Sweep inputs'' button on the ``Waveforms'' tab
applies all eight combinations of $a$, $b$ and $c_{in}$ in turn, and in the
diagram (Fig.~\ref{fig:y-wave}) you can see how $sum$ and $c_{out}$ respond.
Compare with the table in the chapter ``Arithmetic Blocks'': $sum$ equals one when
the number of ones on the inputs is odd, and $c_{out}$ when there are two or three ones.

Now the most interesting part. In Part~I we said that any logic
gate in an FPGA is replaced by a table---a LUT. Let us open the Gowin tab and click
the LUT for $c_{out}$.

\shot{br-adder-lut}{The same adder on Gowin primitives: two LUT3s and the $c_{out}$ table in the inspector}{fig:b-adder-lut}

The synthesizer replaced all five gates with \emph{two} tables (one for each
output). The LUT configuration for $c_{out}$ is the number \code{0xE8}, that is,
\code{11101000$_2$}: the ones are in positions 3, 5, 6 and 7. These are the addresses
\code{011}, \code{101}, \code{110} and \code{111}, which are exactly the combinations
in which at least two of the three input variables equal one.

\begin{ideabox}
The gates from Part~I and the LUT tables are two ways to write one and the same
Boolean function. On paper we choose gates; in an FPGA the synthesizer computes
the table itself. For us the result is the same: a truth table.
\end{ideabox}

\begin{diybox}[The adder and the noise threshold]
\textbf{Task.} Describe a 4-bit adder with carry and a comparison, and check that it computes the same as a chain of full adders in the Designer.

\textbf{In Logra.} The ``From scratch'' project, the ``Code'' tab. Enter \code{assign \{co, s\} = a + b + ci;} with \code{input [3:0] a, b}, \code{input ci}, \code{output [3:0] s}, \code{output co}, and click ``Synthesize''. For the noise threshold, add \code{output loud} and the line \code{assign loud = noise > 4'd9;} (the input \code{noise} is 4 bits wide).

\textbf{Expected result.} The RTL tab shows an adder and a comparator; $9+8+1$ gives \code{s = 2} and \code{co = 1}; \code{loud} equals 0 at \code{noise = 9} and 1 at \code{noise = 10}.

\textbf{How to check.} The ``Waveforms'' tab: set $a=9$, $b=8$, $ci=1$ and compare $s$ and $co$; then sweep the values of \code{noise}. On the Gowin tab write down how many cells the adder took.

\stars{1} repeat: the full adder \code{full\_adder.v} from this chapter, the ``Sweep inputs'' scenario, comparison with the table. \stars{2} modify: make a subtractor \code{assign \{bo, d\} = \{1'b0, a\} - \{1'b0, b\};} and find a pair for which \code{bo = 1} (a borrow). \stars{3} project: the ``Security Panel'' stage---an alarm on a threshold (\code{loud}); acceptance criterion: in the diagram \code{loud} rises exactly when \code{noise} is from 10 to 15.
\end{diybox}

\section{The Decoder: A Shift Instead of Enumeration}

A $2\to 4$ decoder with an enable input can be described either by a table or by
formulas for each output. In Verilog it is easiest to use a shift:
we shift a one left by as many positions as the address says.

\vfile{\codepath/00_theory/decoder2to4.v}{code/00\_theory/decoder2to4.v}

\shot{br-dec-rtl}{The decoder: a shifter and a multiplexer that selects ``all zeros'' when $en = 0$}{fig:b-dec-rtl}

The environment drew a different circuit than ours in the textbook: instead of four AND
gates it showed one \code{<<} shifter and a multiplexer. Both circuits do the same thing.
RTL shows the circuit \emph{by the meaning} of the code, and on the chip it will
be replaced by four LUTs anyway. The diagram in Fig.~\ref{fig:b-dec-wave}
confirms that on any address exactly one line is lit, and at $en = 0$ none.

\shot{br-dec-wave}{Sweeping all 8 input combinations of the decoder}{fig:b-dec-wave}

\begin{diybox}[Zone selection with a decoder]
\textbf{Task.} The security panel watches four zones. The zone address \code{z} (2 bits) must light exactly one of four lamps.

\textbf{In Logra.} The ``Code'' tab: open \code{decoder2to4.v} from this chapter (or enter your own module) and connect the outputs to four zone lamps. Click ``Synthesize''.

\textbf{Expected result.} At $en=1$ one lamp is lit, the one whose number equals the address; at $en=0$ no lamp is lit.

\textbf{How to check.} Waveforms: the ``Sweep inputs'' scenario gives eight combinations, and on each there is at most one one on the output. The RTL tab: find the \code{<<} shifter; compare with your circuit of four AND gates.

\stars{1} repeat: check all 8 rows of the sweep. \stars{2} modify: make a $3\to 8$ decoder using the shift \code{8'd1 << a} and make sure that the circuit became wider while the code is no longer. \stars{3} project: the ``Security Panel'' stage---the alarm zone; acceptance criterion: when the sensor of zone 2 triggers, only lamp 2 is lit, and on reset all lamps go out.
\end{diybox}

\section{The Latch and the Flip-Flop: Two Kinds of Memory Side by Side}

We examined in detail how a latch differs from a flip-flop in the chapter
``Flip-Flops''. Now you can see it in the circuit. Let us describe both elements in one
module: the latch with an \code{always @*} block with a level condition, the flip-flop with
an \code{always @(posedge clk)} block.

\vfile{\codepath/00_theory/dff_vs_latch.v}{code/00\_theory/dff\_vs\_latch.v}

\shot{br-latch-rtl}{On the left, a DFF flip-flop that responds to the edge of the clock signal $C$; on the right, a LATCH with an enable input $EN$}{fig:b-latch-rtl}

The difference in the symbols is the same as in the figures of Part~I:
the flip-flop has a triangle on its clock input (an edge), while the latch's enable
input $EN$ acts by level.

\begin{warnbox}
A latch appears in Verilog \emph{by mistake} more often than by design: if in a
combinational block you do not assign a value under some condition, the synthesizer
inserts a latch to keep the old value. So a \code{LATCH} line
on the RTL diagram that you did not plan should be treated as a
warning signal.
\end{warnbox}

\begin{diybox}[Remember the alarm]
\textbf{Task.} The alarm must stay on even after the door is closed, until a key switches it off. In Part~I a flip-flop did this; describe it in Verilog.

\textbf{In Logra.} The ``Code'' tab, a new module with the inputs \code{clk}, \code{trig}, \code{key} and the output \code{mem}. The module body: \code{always @(posedge clk) if (key) mem <= 1'b0; else if (trig) mem <= 1'b1;}

\textbf{Expected result.} After a pulse on \code{trig} the output \code{mem} stays one until \code{key} arrives; in RTL there is a single DFF flip-flop with reset and set logic.

\textbf{How to check.} Waveforms: apply \code{trig} for one clock, then remove it; \code{mem} must hold; then apply \code{key}. There must be no LATCH elements in RTL.

\stars{1} repeat: find the DFF and the LATCH in \code{dff\_vs\_latch.v}. \stars{2} find the bug: remove the \code{else} branch in the \code{always @*} block that sets \code{y} only when \code{en = 1}, and explain where the latch in the circuit came from. \stars{3} project: the ``Security Panel'' stage---remembering the alarm; acceptance criterion: in the diagram \code{mem} does not release the alarm after the door closes and is cleared only by \code{key}.
\end{diybox}

\section{The Counter: A Register Plus an Adder}

The synchronous counter from the chapter ``Registers and Counters'' consists of a register and
a circuit that adds one. Verilog describes this with a single line
\code{q <= q + 1}, and reset and count enable with conditional statements.

\vfile{\codepath/00_theory/counter4.v}{code/00\_theory/counter4.v}

\shot{br-cnt-rtl}{The RTL diagram of a counter: a $+1$ adder and a 4-bit register with synchronous reset and an enable input}{fig:b-cnt-rtl}

In the diagram it is easy to find all the parts we talked about: the adder
(a circle with a plus), the register (a DFF rectangle), the enable input \code{EN},
the synchronous reset \code{SRST}. The register's output returns to the adder's
input---this is the feedback that turns a combinational
adder into a counter.

The diagram (Fig.~\ref{fig:b-cnt-wave}) shows the count. First the
reset is active (\code{rst = 1}), then counting goes on the rising clock edges: $0, 1,
2, \ldots, 15$, and then $q$ returns to zero---the counter's modulus is 16.

\shot{br-cnt-wave}{The counter at work: reset, counting and overflow}{fig:b-cnt-wave}

\begin{diybox}[An alarm duration counter]
\textbf{Task.} The security panel must know how many clocks the alarm lasts, and must not overflow: on reaching 15, the counter stops.

\textbf{In Logra.} The ``Code'' tab. Take \code{counter4.v} and replace the counting condition with \code{if (en \&\& q != 4'd15) q <= q + 4'd1;}. Synthesize and connect \code{q} to four LEDs.

\textbf{Expected result.} At \code{en = 1} the value grows $0,1,\dots,15$ and stays at 15 (saturation); in the original \code{counter4} it would be 0 after 15.

\textbf{How to check.} Waveforms: set the number of cycles to 20 and keep \code{en = 1}. In RTL find the ``+1'' adder and the register with an enable input, as in Fig.~\ref{fig:b-cnt-rtl}.

\stars{1} repeat: run the original \code{counter4} and read the modulus 16 from the diagram. \stars{2} modify: make a modulo-10 counter (0 again after 9). \stars{3} project: the ``Security Panel'' stage---alarm time; acceptance criterion: while \code{state} equals ALARM, the counter grows to 15 and stays at 15, and on reset returns to 0.
\end{diybox}

\section{The Finite-State Machine: A Traffic Light}

The finite-state machine from Chapter~\ref{ch:sync} is a
state register and transition logic. A Moore machine for a traffic light has three
states. In each of them one lamp is lit, and the transition is determined only by
the number of elapsed clocks. We are not meeting state machines again: you have
already built the circuit and the graph from blocks, and here the same structure is written as code.

\begin{figure}[H]
\centering
\begin{tikzpicture}[st/.style={circle, draw=FpgaBlue, fill=FpgaLight, minimum size=18mm, font=\sffamily\small, align=center, thick}]
  \node[st] (g) at (0,0) {GREEN\\{\scriptsize 4 clocks}};
  \node[st] (y) at (5,0) {YELLOW\\{\scriptsize 2 clocks}};
  \node[st] (r) at (10,0) {RED\\{\scriptsize 4 clocks}};
  \draw[arr] (g) -- node[above, lbl] {$t=3$} (y);
  \draw[arr] (y) -- node[above, lbl] {$t=1$} (r);
  \draw[arr] (r) to[bend left=35] node[below, lbl] {$t=3$} (g);
  \draw[arr] (-2.4,0) -- node[above, lbl] {rst} (g);
\end{tikzpicture}
\caption{The traffic light state graph. $t$ is the counter of clocks spent in the state}
\label{fig:b-fsm-graph}
\end{figure}

\vfile{\codepath/00_theory/traffic_fsm.v}{code/00\_theory/traffic\_fsm.v}

In the code the three parts of the machine are clearly visible: the state register
(\code{state}), the transition logic (the \code{case} statement in the block on the clock
edge) and the output logic (the second \code{case} block, depending only on
\code{state}). In the diagram (Fig.~\ref{fig:b-fsm-wave}) you can see what
a Moore machine is supposed to do: the output \code{lamp} changes exactly when
\code{state} changes.

\shot{br-fsm-wave}{The traffic light: in each state of the machine one lamp is lit ($0x1$ is green, $0x2$ is yellow, $0x4$ is red)}{fig:b-fsm-wave}

\begin{diybox}[The ``Guard'' machine: three states]
\textbf{Task.} Following the traffic light example, build a security panel machine with the states OFF, ARMED, ALARM: the button \code{arm} takes OFF to ARMED, a sensor trigger takes ARMED to ALARM, and reset takes it to OFF.

\textbf{In Logra.} The ``Code'' tab. First run \code{traffic\_fsm.v}, then on its basis write a machine with a two-bit \code{state} and three \code{case} branches. Do not forget the \code{default} branch.

\textbf{Expected result.} The traffic light holds the lamps for 4, 2 and 4 clocks in the order green, yellow, red. After \code{arm}, the guard machine stays in the ARMED state until a sensor triggers.

\textbf{How to check.} Waveforms: for the traffic light, \code{lamp} takes the values \code{0x1}, \code{0x2}, \code{0x4} and changes simultaneously with \code{state}. For the panel, apply \code{arm}, then \code{door}, and make sure that \code{state} went through 0, 1, 2. An extra transition you will see at once.

\stars{1} repeat: recount the lamp durations of the traffic light from the diagram. \stars{2} modify: increase green to 6 clocks (the constant in the comparison with \code{t}) and check the diagram. \stars{3} project: the ``Security Panel'' stage---the state machine; acceptance criterion: the three-state graph is drawn, the diagram shows all three states and not a single unintended transition.
\end{diybox}

\section{The Mealy Machine: A Vending Machine}

In Chapter~\ref{ch:fsmdesign} we drew the graph of a vending machine and checked
it by hand against a table. Let us write the same machine in Verilog. It has three parts:
the state register, the transition logic and the output logic, and the output depends on both the
state and the current coin.

\vfile{\codepath/00_theory/vending.v}{code/00\_theory/vending.v}

Three techniques protect the description from typical mistakes: the default value
\code{next = state} (so that no latch appears), the \code{default} branch (so that the
machine does not hang in an unused code) and synchronous reset.


Let us check the machine with the same coins as in the table of Chapter~\ref{ch:fsmdesign}.
The ``Waveforms'' tab, the ``Events'' section, lets you set the coin pulses by
periods (Fig.~\ref{fig:n-vend-wave}).

\shot{n-vend-wave}{The Mealy machine: \code{vend} rises in the same clock in which the last coin arrives}{fig:n-vend-wave}

The signal \code{vend} rises \emph{within} the clock cycle, as soon as the ``deciding'' coin
appears on the input, without waiting for the next edge. This is the
difference between a Mealy machine and a Moore machine. If you need a ``clean'' output, it
is passed through a flip-flop, and it lags by one clock.

\begin{diybox}[Entering the code 1-0-1 as a Mealy machine]
\textbf{Task.} The panel lock is released if the bits 1, 0, 1 are entered in a row. Make a sequence detector whose output \code{unlock} rises in the same clock in which the last bit arrives.

\textbf{In Logra.} The ``Code'' tab. The model is \code{vending.v}: the same three parts (state register, \code{next}, output). There are three states: ``nothing'', ``1'', ``10''. Or start with the lesson ``A 101 sequence detector'' on the ``Learn'' tab.

\textbf{Expected result.} The input series \code{1 0 1} raises \code{unlock}; the series \code{1 1 0} and \code{0 1 1} do not; the series \code{1 0 1 0 1} gives two pulses (overlap).

\textbf{How to check.} Waveforms, the ``Events'' section: set the \code{code\_bit} pulses by periods, as in Fig.~\ref{fig:n-vend-wave}. The output changes within the clock, not after it: this is what a Mealy machine looks like.

\stars{1} repeat: run a sequence of coins through \code{vending.v} and find the clock where \code{vend = 1}. \stars{2} modify: make the code 1-1-0 and check that 1-0-1 no longer works. \stars{3} project: the ``Security Panel'' stage---code entry; acceptance criterion: only the exact series 1-0-1 returns the panel to the OFF state.
\end{diybox}

\section{The FIFO Queue}

The queue from Chapter~\ref{ch:fsmdesign} consists of an array of words, two
pointers and a word counter. The module \code{fifo8} stores eight words of
eight bits each.

\vfile{\codepath/00_theory/fifo8.v}{code/00\_theory/fifo8.v}

The synthesizer will place the array \code{mem} in the chip's memory, and the pointers are
ordinary counters that ``wrap around'' by themselves after seven. Writing and reading are
allowed only while the queue is not full and not empty.

\shot{n-fifo-wave}{The queue at work: the words \code{0B}, \code{16}, \code{21} are written, then read in the same order}{fig:n-fifo-wave}

In the diagram \code{count} grows with each write and falls with each read,
\code{empty} is cleared after the first write, and the words come out on \code{dout} in
the order they were written.

\begin{diybox}[An alarm log in a queue]
\textbf{Task.} The panel must remember the numbers of the last zones where the alarm went off, and output them on request in the same order.

\textbf{In Logra.} The ``Code'' tab: open \code{fifo8.v}, synthesize. Apply the zone numbers to \code{din} (on the 8-bit bus the low two bits are enough) and the pulses \code{wr}.

\textbf{Expected result.} The written numbers come out on \code{dout} in the order of writing; \code{count} grows on writing and falls on reading, \code{empty} goes low after the first write.

\textbf{How to check.} Waveforms: write three words and read them, as in Fig.~\ref{fig:n-fifo-wave}. Then write 9 words in a row: the ninth must not be accepted (\code{full = 1}). The array \code{mem} on the Gowin tab will take a memory block or cells.

\stars{1} repeat: the example ``Four-word FIFO'' (\code{fifo-four}) from the ``From an example'' card. \stars{2} modify: reduce the depth to four words (pointers of 2 bits, \code{count} up to 4). \stars{3} project: the ``Security Panel'' stage---the log; acceptance criterion: after three triggerings of zones 1, 3, 2 reading returns exactly 1, 3, 2.
\end{diybox}

\section{The Hamming Code}

The encoder, the channel with an error and the decoder from Chapter~\ref{ch:codes} consist entirely of
exclusive ORs. For testing, the module has an input \code{err}: a mask
that flips the chosen bits in the ``channel''.

\vfile{\codepath/00_theory/hamming74.v}{code/00\_theory/hamming74.v}


The checking module \code{tb\_hamming74.v} (a testbench, Chapter~\ref{ch:sim}; it lies
alongside) goes through all 16
data values and all seven single errors and makes sure that the data is recovered every
time and that the syndrome points to the damaged position.

\begin{diybox}[Protecting a zone code from an error]
\textbf{Task.} A zone number (4 bits) travels over a long wire, and one bit in it may flip. Check that the Hamming code corrects any single error.

\textbf{In Logra.} The ``Code'' tab: \code{hamming74.v}. The input \code{d} is your zone number, \code{err} is the damage mask. The testbench \code{tb\_hamming74.v} lies alongside.

\textbf{Expected result.} At \code{err = 0000100} the outputs \code{syn} equal 3, \code{q} equals the original \code{d}. At \code{err = 0000000} the syndrome is 0.

\textbf{How to check.} Waveforms: set \code{d = 1011} and in turn \code{err} with a single one in each of the seven positions. In RTL find the XOR chains (the encoder and the syndrome computation).

\stars{1} repeat: go through all seven positions and write down the syndromes. \stars{2} find the bug: apply \code{err} with two ones, make sure that \code{q} is now wrong, and explain why the code does not cure such a case. \stars{3} project: the ``Security Panel'' stage---code protection; acceptance criterion: for all 16 zone numbers and all seven single errors \code{q} matches \code{d}.
\end{diybox}

\section{A Three-State Pin}

In Chapter~\ref{ch:electrical} we found out that an output can be
\emph{disconnected} (the state $Z$). In Verilog this is written with a conditional
assignment that outputs the value \code{z} in the disabled state.
A bidirectional pin (\code{inout}) adds to this the reading of what
is actually on the wire.

\vfile{\codepath/00_theory/io_pad.v}{code/00\_theory/io\_pad.v}

\shot{n-iopad-rtl}{A bidirectional pin: a tri-state buffer on the output and an ordinary reading of the level on the input}{fig:n-iopad}

\begin{diybox}[A siren on a shared line]
\textbf{Task.} The siren and the panel indicator share one pin: the line belongs to the siren while the alarm lasts, and the rest of the time it is released ($Z$) and can accept a signal from outside.

\textbf{In Logra.} The ``Code'' tab: open \code{io\_pad.v}, synthesize. Connect \code{oe} to the ``alarm'' flag.

\textbf{Expected result.} At \code{oe = 1} the pin \code{pad} shows \code{to\_pad}; at \code{oe = 0} the pin is not driven by us, and \code{from\_pad} shows what is applied from outside.

\textbf{How to check.} RTL: a tri-state buffer on the output and an ordinary reading of the level (Fig.~\ref{fig:n-iopad}). Waveforms: toggle \code{oe} and watch \code{from\_pad}. Pin assignment is on the ``Pins'' tab.

\stars{1} repeat: go through both modes of \code{oe}. \stars{2} modify: split the pin into two directions by bits (two bits output, two input). \stars{3} project: the ``Security Panel'' stage---the siren on a shared line; acceptance criterion: outside an alarm \code{oe = 0} and the line is free.
\end{diybox}

\section{A DAC Based on PWM}

The module from Chapter~\ref{ch:analog}: a counter that runs in a circle, and
a comparison with the required value.

\vfile{\codepath/00_theory/pwm_dac.v}{code/00\_theory/pwm\_dac.v}

The counter \code{cnt} runs from 0 to 255, and the output equals one while
\code{cnt < level}. At \code{level = 64} the one is held for 64 clocks out of 256.

\shot{n-pwm-wave}{A PWM DAC with the code 64: the output equals one for the first quarter of the counter period}{fig:n-pwm}

\begin{diybox}[Buzzer volume through PWM]
\textbf{Task.} The panel buzzer's volume is set by the number \code{level}. Get the required fraction of time with a one on the output.

\textbf{In Logra.} The ``Code'' tab: \code{pwm\_dac.v}. Also open the example ``Brightness with PWM'' (the ``From an example'' card) and compare.

\textbf{Expected result.} At \code{level = 64} the output equals 1 for a quarter of the counter period (25\,\%), at \code{level = 192} for three quarters (75\,\%), at 0 it is always 0.

\textbf{How to check.} Waveforms: set \code{level}, set the number of cycles to 512 (two periods) and count how many clocks \code{pwm = 1}. The average voltage equals $D\cdot 3.3$~V: for $D=25\,\%$ you get about 0.8~V.

\stars{1} repeat: \code{level = 64}, a period of 256 clocks, one for the first 64. \stars{2} modify: make the counter 6 bits wide and recompute the fractions. \stars{3} project: the ``Security Panel'' stage---the buzzer; acceptance criterion: in the ALARM state the output \code{pwm} has a duty cycle of 50\,\%, outside an alarm it is 0.
\end{diybox}

\section{A Successive-Approximation ADC}

The state machine from Chapter~\ref{ch:analog} uses the comparator signal \code{cmp} to find
the code by bisection. The file has two modules: the machine itself \code{sar\_adc} and
the teaching wrapper \code{sar\_demo}, where the ``input voltage'' is set by the constant
179, and the comparator is built from a comparison operator.

\vfile{\codepath/00_theory/sar_adc.v}{code/00\_theory/sar\_adc.v}

\shot{n-sar-wave}{The search for the number 179: the DAC code passes through the values \code{80}, \code{C0}, \code{A0}, \code{B0}, \code{B8}, \code{B4}, \code{B2}, \code{B3}}{fig:n-sar}

In the diagram you can clearly see the binary search that we examined in the table:
on each clock \code{cmp} decides whether the bit under test stays one, and
after eight clocks \code{result} becomes \code{0xB3}, that is, 179.

\begin{diybox}[A noise sensor through an ADC]
\textbf{Task.} The noise level in the panel arrives from a comparator, but you need a number. Find the code by bisection in 8 clocks.

\textbf{In Logra.} The ``Code'' tab: \code{sar\_adc.v}, top-level module \code{sar\_demo}. The ``input voltage'' is set by the constant 179.

\textbf{Expected result.} The DAC code passes through the values \code{0x80}, \code{0xC0}, \code{0xA0}, \code{0xB0}, \code{0xB8}, \code{0xB4}, \code{0xB2}, \code{0xB3}; after 8 clocks \code{done = 1} and \code{result = 0xB3}.

\textbf{How to check.} Waveforms: the signals \code{dac}, \code{cmp}, \code{result}, \code{done}, the number of cycles 20, a \code{start} pulse. Figure~\ref{fig:n-sar} shows the expected picture.

\stars{1} repeat: run \code{sar\_demo} and check all eight steps. \stars{2} modify: set the input value to 100 and predict the steps in advance, \code{result} must become \code{0x64}. \stars{3} project: the ``Security Panel'' stage---noise through the ADC; acceptance criterion: the comparison \code{result[7:4] > 4'd9} gives \code{loud = 1} exactly for results from \code{0xA0} and above (for 179 this is one, for 100 zero).
\end{diybox}

\section{The Security Panel in Verilog}

Here are all the blocks together. In Part~I the panel grew in stages: a door sensor,
arming and night mode, the ``2 out of 3'' vote, the noise threshold, remembering the alarm,
the state machine and code entry. Now each stage corresponds to a few lines of
Verilog. The module introduces nothing new: it repeats the circuit you assembled
from blocks.

\begin{vcode}
module guard_panel (
    input        clk, rst,                     // clock and synchronous reset
    input        door, window, motion,         // sensors, 1 = triggered
    input        night,                        // night mode
    input  [3:0] noise,                        // noise level
    input        arm_btn,                      // the "arm" button
    input        code_bit, code_step,          // code entry: the bit and the "accept" pulse
    output       alarm_led, armed_led,
    output       siren,
    output reg [1:0] state
);
    localparam OFF = 2'd0, ARMED = 2'd1, ALARM = 2'd2;

    // Stages 2, 3, 6: the combinational part
    wire vote     = (door & window) | (door & motion) | (window & motion);
    wire loud     = noise > 4'd9;
    wire intruder = vote | loud | (door & night);

\end{vcode}

The second half of the module: code entry, the state machine and the siren.

\begin{vcode}
    // Stage 9: entering the code 1-0-1 and the button edge
    reg [2:0] hist;
    reg       arm_d;
    always @(posedge clk) begin
        arm_d <= arm_btn;
        if (rst)            hist <= 3'b000;
        else if (code_step) hist <= {hist[1:0], code_bit};
    end
    wire arm_pulse = arm_btn & ~arm_d;
    wire unlock    = code_step & ({hist[1:0], code_bit} == 3'b101);

    // Stages 7 and 9: the state machine OFF -> ARMED -> ALARM (ALARM is remembered)
    always @(posedge clk) begin
        if (rst) state <= OFF;
        else case (state)
            OFF:     if (arm_pulse) state <= ARMED;
            ARMED:   if (unlock)        state <= OFF;
                     else if (intruder) state <= ALARM;
            ALARM:   if (unlock) state <= OFF;
            default: state <= OFF;
        endcase
    end

    // Stage 8: the counter for the beeping (tone and pause)
    reg [23:0] cnt;
    always @(posedge clk) cnt <= rst ? 24'd0 : cnt + 24'd1;

    assign alarm_led = (state == ALARM);
    assign armed_led = (state == ARMED);
    assign siren     = alarm_led & cnt[22] & cnt[13];
endmodule
\end{vcode}

Let us see where everything came from. The lines with \code{wire} are the combinational part;
in RTL this is an adder-comparator and a few gates. The
\code{always @(posedge clk)} block with \code{hist} is a shift register
of three flip-flops that catches the code 1-0-1. The \code{case} statement is a machine of
three states; the ALARM state is held until the correct code arrives: this is
the alarm memory. The counter \code{cnt} and two of its bits produce a beep with pauses.
At a frequency of 27~MHz the bit \code{cnt[13]} has a frequency of about 1.6~kHz (the tone), and
\code{cnt[22]} about 3.2~Hz (the pause between beeps).

\begin{warnbox}
The signals of buttons and sensors come from outside and are in no way related to the panel's
clock. Before use they are passed through two flip-flops and a contact-bounce
filter (Chapter~\ref{ch:clocks}). In the Logra simulation the signals are ideal, but on a board
without a synchronizer the panel will occasionally ``make mistakes''.
\end{warnbox}

\begin{diybox}[The Security Panel: the summary of Part I in Verilog]
\textbf{Task.} Build \code{guard\_panel} from this chapter in Logra, check all the stages and prepare it for the board.

\textbf{In Logra.} The ``From scratch'' project, the ``Code'' tab: paste the module, click ``Synthesize''. For a demonstration in Waveforms replace \code{cnt[22]} and \code{cnt[13]} with \code{cnt[3]} and \code{cnt[1]}, otherwise you would wait millions of clocks for the beep.

\textbf{Expected result.} After \code{arm\_btn} the panel is in the ARMED state. Triggering of two sensors out of three, noise above 9, or the door at night take it to ALARM, and the siren sounds. Closing the door does not cancel the alarm. The series \code{1, 0, 1} on \code{code\_bit} with \code{code\_step} pulses returns OFF.

\textbf{How to check.} Waveforms: the signals \code{state}, \code{door}, \code{window}, \code{noise}, \code{alarm\_led}, \code{siren}; the number of cycles 120. RTL: find the comparator, the three \code{hist} flip-flops, the state machine. The Gowin tab will show the number of cells. Then on the ``Pins'' tab assign the buttons, LEDs and the buzzer and save the project.

\stars{3} Final acceptance (the ``Security Panel'' stage, board): (1) the panel arms on the button and does not react to sensors in the OFF state; (2) any of the three causes takes it to ALARM, closing the door does not cancel the alarm; (3) the alarm can be cleared only with the code 1-0-1, any other series does not work; (4) on the Tang Nano 20K board the buttons are connected through a synchronizer and debounce, and the buzzer gives a beep with pauses.
\end{diybox}

\begin{ideabox}
The whole ``path from idea to circuit'' repeats every time:
a statement of behavior $\to$ a Verilog description $\to$ a circuit of gates and
flip-flops $\to$ the chip's LUT tables and flip-flops. The ability to move between
these levels is the main skill of a digital designer.
\end{ideabox}

\begin{taskbox}[Tasks]
\begin{enumerate}
  \item Describe in Verilog a 4-bit adder with a carry input and an overflow output
        and check it on the ``Waveforms'' tab. How many ALU elements
        does the synthesizer use on the Gowin tab?
  \item Add an input \code{up} to \code{counter4}: at \code{up = 0} the counter must
        count down. What changed in the RTL diagram?
  \item Change the traffic light so that before the green light the yellow
        lights up together with the red for two clocks. How many states are needed now?
  \item Describe a $4:1$ multiplexer with the statement \code{assign y = d[s]} and
        compare its circuit on the RTL tab with the figure in the chapter ``Multiplexers and Demultiplexers''.
\end{enumerate}
\end{taskbox}
